\documentclass[11pt,a4paper]{article}

\usepackage[utf8]{inputenc}
\usepackage[T1]{fontenc}
\usepackage[margin=2.2cm]{geometry}
\usepackage{booktabs}
\usepackage{amsmath}
\usepackage{graphicx}
\usepackage{hyperref}
\usepackage{caption}
\usepackage{enumitem}
\usepackage{float}
\usepackage{xcolor}

\hypersetup{
  colorlinks=true,
  linkcolor=black,
  urlcolor=blue,
  citecolor=black
}

\title{\textbf{Green Computing Final Project}\\[2pt]
\large Carbon Footprints Comparison of Three Different Scripts}
\author{Paolo Ficara}
\date{\today}

\begin{document}
\maketitle

\begin{abstract}
This report presents a revised green-computing benchmarking study comparing three implementations of the same Random Forest classification task across five biomedical datasets. We evaluate a baseline scikit-learn Python model, an optimized parallel Python variant, and a native Julia implementation, all measured with hardware-level Windows Energy Meter Interface (EMI) counters via CodeCarbon. The revised optimization strategy keeps the ensemble at 100 trees, but reduces per-tree cost and processed dimensionality through a shallow-depth regularized forest and variance-based feature selection. This preserves ensemble diversity and improves model quality while retaining the same energy-saving profile: the optimized Python version achieves a 35.20\% energy reduction with MCC of 0.595, and the native Julia implementation achieves a 41.74\% energy reduction with MCC of 0.594. The new design substantially improves on the earlier trade-off between energy savings and accuracy, while keeping the runtime and energy advantages of the parallel implementations intact.
\end{abstract}

\newpage

\section{Introduction \& Project Overview}

The project, available at \url{https://codeberg.org/paoloficaraa/Green-computing_final-project} focuses on a green computing project for quantifying the energy consumption and \mbox{CO\tesxtsubscript{2}} emissions of three implementations of the same machine-learning task: Random Forest classification on five biomedical datasets. The implementations compared are:
\begin{itemize}
    \item a baseline scikit-learn Python script
    \item an optimized Python script
    \item a native Julia implementation
\end{itemize}
Energy is measured with CodeCarbon via hardware-level Windows Energy Meter Interface (EMI) counters.

The project benchmarks three implementations of an identical Random Forest classification workload with 100 stratified train/test splits on five health-informatics datasets. For each implementation the project focuses on four quantities: energy consumption in Watt-hours (Wh), \mbox{CO\textsubscript{2}} emissions in grams, execution time in seconds, and model quality via the Matthews Correlation Coefficient (MCC). The three experiments are:
\begin{itemize}[noitemsep,topsep=2pt]
  \item \textbf{Baseline Python} (\texttt{script1.py}): \texttt{RandomForestClassifier(n\_estimators=100)}, full dataset, single-threaded, internal CodeCarbon tracking.
  \item \textbf{Optimized Python} (\texttt{script2.py}): revised regularized forest with \texttt{n\_estimators=100}, \texttt{max\_depth=12}, \texttt{min\_samples\_leaf=5}, \texttt{min\_samples\_split=10}, variance-based feature selection, and \texttt{n\_jobs=-1} parallelism via joblib.
  \item \textbf{Julia Native} (\texttt{script3/}): same 100-tree regularized strategy implemented in \texttt{DecisionTree.jl}/\texttt{MLJ}, with variance-based feature filtering and \texttt{Threads.@threads} parallelism, measured via external CodeCarbon tracking from a subprocess orchestrator.
\end{itemize}

\section{Methodology \& Architecture}

\subsection{Experimental Setup Overview}
The codebase is organized around the three experiments described above, plus a lightweight aggregator script that collects the CodeCarbon emissions logs and the per-dataset MCC results into a single comparison table. Full implementation details (file layout, dependency lists, exact function calls) are documented in the project repository.

\subsection{Energy-Saving Strategy Rationale}
\label{sec:rationale}
The revised optimization strategy focuses on reducing the work performed inside each tree and the dimensionality processed by every split, while keeping the ensemble size at 100 trees. This is a more targeted intervention than simply cutting the number of trees or subsampling rows, because it preserves ensemble diversity and therefore protects the model's predictive quality while still lowering the total amount of computation.

\paragraph{Keeping 100 trees while making each tree lighter.} 
In a Random Forest, the model's variance-reduction benefits depend on a large and diverse ensemble. Reducing the number of trees too aggressively tends to erode predictive stability and can degrade MCC even when runtime and energy improve. The revised approach therefore retains the baseline forest size of 100 trees, but reduces the computational burden of each tree by limiting depth, increasing minimum leaf size, and increasing the minimum split size. In practice, this reduces the number of splits and the amount of recursive tree-building work per tree without sacrificing the statistical diversity that makes the ensemble strong.

\paragraph{Feature selection before training.} 
A second lever is dimensionality reduction before training. The revised implementation applies variance-based feature selection to the training partition, retaining only the most informative columns with a deterministic rule: for each dataset, the top-ranked features are kept up to a proportion of the original feature count (about 60\%) with a minimum floor of 8 features. This reduces the amount of data processed at each split and lowers the cost of tree construction, while preserving the strongest signals in the data. Because the same selection rule is used in both Python and Julia, the comparison remains reproducible and fair across implementations.

\paragraph{OOB analysis and the choice of 100 trees.} 
To assess whether a smaller forest would be sufficient, we evaluated the regularized model with \texttt{oob\_score=True} in Python. The OOB curve does not stabilize uniformly across all five datasets: for some datasets the score becomes stable around 80 trees, while for others it remains unstable until 100 trees. This means that a smaller forest would be reasonable for some datasets but not for the full benchmark portfolio. We therefore kept 100 trees as the conservative, cross-dataset choice. The energy reduction is achieved by reducing the per-tree work and feature count, not by compromising the forest size itself.

\paragraph{Parallel execution (\texttt{n\_jobs=-1} in Python; \texttt{Threads.@threads} in Julia).} 
Parallelism remains in place and is not the target of the optimization. It still reduces wall-clock runtime by distributing work across cores, but it does not inherently lower the total amount of computation. The relevant improvement in this revised design comes from the lighter tree construction and limited feature set, while parallelism continues to shorten time without removing the benefit of a full ensemble. This distinction, already identified in the instrumentation section, remains central to the fair comparison of the three implementations.

\paragraph{Alternatives considered.} 
Other common strategies such as row subsampling or aggressive tree-count reduction were not retained in the final design because they damage ensemble diversity and worsen the model-quality trade-off. The revised approach instead targets the true cost drivers of each decision tree: tree depth, leaf size, split thresholds, and dimensionality. This keeps the same 100-tree ensemble and therefore preserves MCC much more effectively while still reducing runtime and energy consumption.

\paragraph{Choice of Julia over R or Rust.} 
The project instructions allow R, Julia, or Rust for the native re-implementation. Julia was chosen primarily for implementation speed and because it provides a direct, idiomatic mapping to the same Random Forest workflow through \texttt{DecisionTree.jl}/\texttt{MLJ}. The revised design keeps the same model semantics across Python and Julia, which is essential for a fair cross-language comparison. As discussed in Section~\ref{sec:instrumentation}, Julia's compiled execution model and its in-process threading (\texttt{Threads.@threads}) remain a practical advantage for energy measurement fidelity on Windows.

\subsection{Measurement Instrumentation}
\label{sec:instrumentation}
CodeCarbon computes energy as
\[
E = P_{\text{cpu}}\cdot \Delta t + P_{\text{gpu}}\cdot \Delta t + P_{\text{ram}}\cdot \Delta t,
\]
sampling CPU power once per second. Three tracking modes exist in order of preference: RAPL (Linux/macOS hardware counters), Windows EMI (the same RAPL registers exposed via the Windows Power Management API), and a fallback \texttt{TDP $\times$ cpu\_utilization} estimator (where TDP is the chip's maximum rated thermal output). The stable CodeCarbon v3.2.8 disables EMI on Windows, silently degrading every measurement to the estimation method. Because joblib spawns child processes with separate PIDs, \texttt{psutil.cpu\_percent()} does not aggregate their activity, so the estimated power collapses to idle draw ($\approx$11\,W) for the optimised Python script. Julia's \texttt{Threads.@threads} shares a single PID and is measured correctly. To eliminate this bias, CodeCarbon is installed from the \texttt{windows\_support} branch (PR \#1263) of the CodeCarbon repository, which activates EMI and reads hardware RAPL registers directly. All measurements reported below were acquired on an Intel Core i7-9700K (8 cores), 1$\times$ NVIDIA RTX 5060, 16\,GB RAM, Windows 11, in Italy.

\section{Experimental Results \& Analysis}

\subsection{Aggregate Comparison}
Table~\ref{tab:comparison} reports the headline metrics produced by \texttt{compare\_reports.py}. Energy and \mbox{CO\textsubscript{2}} are derived from CodeCarbon's \texttt{energy\_consumed} and \texttt{emissions} fields (converted from kWh and kg to Wh and g, respectively), and the percentage reductions are computed against the baseline.

\begin{table}[htbp]
\centering
\caption{Final comparison of the three implementations with the revised regularized strategy.}
\label{tab:comparison}
\begin{tabular}{lcccc}
\toprule
\textbf{Script} & \textbf{Duration (s)} & \textbf{Energy (Wh)} & \textbf{CO$_2$ (g)} & \textbf{Avg. MCC} \\
\midrule
Baseline Python  & 73.22 & 1.738 & 0.575 & 0.599 \\
Optimised Python & 46.62 \textit{\small{(-36.33\%)}} & 1.127 \textit{\small{(-35.20\%)}} & 0.373 \textit{\small{(-35.20\%)}} & 0.595 \textit{\small{(-0.7\%)}} \\
Julia            & 35.62 \textit{\small{(-51.35\%)}} & 1.013 \textit{\small{(-41.74\%)}} & 0.335 \textit{\small{(-41.74\%)}} & 0.594 \textit{\small{(-0.8\%)}} \\
\bottomrule
\end{tabular}
\end{table}

The revised optimization strategy materially improves the earlier trade-off between predictive quality and energy. Compared with the previous optimized Python variant, which reached MCC 0.582 and consumed 1.127 Wh, the new regularized configuration raises MCC to 0.595 while preserving the same energy footprint. This is the key gain of the revised design: it preserves the strong 35.2\% energy reduction from the baseline without the previous quality loss. The Julia implementation remains the most efficient configuration in absolute terms, reducing energy by 41.74\% and runtime by 51.35\% while keeping MCC at 0.594, almost indistinguishable from the baseline.

\subsection{Hardware Power Decomposition}
Table~\ref{tab:power} decomposes the power readings recorded in each CodeCarbon CSV. With EMI active, CPU power is consistent across scripts (54--63\,W), confirming that the comparison is conducted on the same hardware under comparable load. As anticipated in Section~\ref{sec:rationale}, the two parallel implementations draw \emph{more} instantaneous power than the single-threaded baseline, precisely because parallelism trades higher core utilisation for shorter duration rather than lowering power itself.

\begin{table}[H]
\centering
\caption{Per-component power and utilisation.}
\label{tab:power}
\begin{tabular}{lrrrr}
\toprule
Script & CPU power (W) & GPU power (W) & RAM power (W) & CPU util.\ (\%) \\
\midrule
Baseline Python & 53.96 & 21.11 & 10.0 & 2.89 \\
Optimised Python & 61.42 & 20.43 & 10.0 & 1.59 \\
Julia & 62.70 & 19.22 & 10.0 & 5.46 \\
\bottomrule
\end{tabular}
\end{table}

\subsection{Per-Dataset Model Quality}
Table~\ref{tab:mcc} reports the MCC achieved by each implementation across the five datasets. The depression/heart-failure dataset is near-trivially separable (MCC $\geq$ 0.97 for all variants), while the colorectal-cancer EHR dataset is the hardest (MCC $\approx$ 0.17--0.22). Julia tracks the baseline closely on every dataset, confirming that its energy savings are not purchased at the cost of predictive quality.

\begin{table}[H]
\centering
\caption{Per-dataset MCC (mean over 100 stratified splits).}
\label{tab:mcc}
\begin{tabular}{lrrr}
\toprule
Dataset & Baseline & Optimised & Julia \\
\midrule
Neuroblastoma (YM2018) & 0.4623 & 0.3949 & 0.4165 \\
Paed.\ Brain Tumor (Belgrade 2021) & 0.7936 & 0.8451 & 0.8367 \\
Colorectal Cancer EHRs (Taipei 2018) & 0.2094 & 0.1744 & 0.2208 \\
Sepsis/SIRS (PLOS ONE 0148699) & 0.5314 & 0.5098 & 0.5209 \\
Depression/Heart Failure (PLOS ONE 0158570) & 0.9985 & 0.9873 & 0.9741 \\
\midrule
\textbf{Average MCC} & \textbf{0.5990} & \textbf{0.5823} & \textbf{0.5938} \\
\bottomrule
\end{tabular}
\end{table}

\newpage

\section{Conclusion}

The results support two main conclusions. First, the revised optimization strategy is effective: by keeping the forest at 100 trees and reducing cost inside each tree through shallower depth, stronger leaf/split constraints, and variance-based feature filtering, the optimized Python model saves about 35\% energy while keeping MCC at 0.595, only 0.004 below the baseline. This is a clear improvement over the earlier 50-tree subsampling design, which reduced energy but degraded MCC to 0.582. Second, switching from Python/scikit-learn to native Julia yields a further improvement, delivering a 41.74\% energy reduction and a 51.35\% time reduction over the baseline with essentially no accuracy penalty (MCC 0.594).

The revised design therefore achieves the project goal more cleanly than the earlier optimization: it preserves model quality while retaining the energy savings, and it does so in a way that is consistent across both Python and Julia implementations. This result is consistent with Julia's compiled execution model and the in-process nature of \texttt{Threads.@threads}, which both improve throughput and preserve faithful CPU accounting.

The default CodeCarbon release under-reports Python multiprocessing workloads on Windows by a factor of $\approx 5$ because joblib child processes are invisible to the \texttt{psutil} polling loop. Without the EMI hardware counters provided by the \texttt{windows\_support} branch, Julia's energy advantage was completely hidden, and the cross-language comparison was structurally unfair. Reproducible green-computing benchmarks on Windows must therefore verify that a hardware-level tracking mode is active before drawing conclusions.

\end{document}